\documentclass[11pt]{article}
\usepackage[margin=1in]{geometry}
\usepackage{times}
\usepackage{hyperref}
\hypersetup{colorlinks=true, linkcolor=black, urlcolor=blue, citecolor=black}
\title{One Man and a Block and Tackle: Coral Castle and the Antigravity Attribution}
\author{Flyxion \\ \small Independent Researcher}
\date{}
\begin{document}
\maketitle

\begin{abstract}
Edward Leedskalnin single-handedly quarried, carved, and assembled Coral Castle in Florida between roughly 1923 and his death in 1951, moving and precisely placing coral limestone blocks weighing up to thirty tons using methods he never demonstrated publicly and did not fully document, giving rise to popular claims that he had discovered a form of magnetic or antigravity levitation technique connected to lost knowledge of the pyramid builders. This paper argues that the absence of a witnessed construction method is fully explained by Leedskalnin's own well-documented preference for working alone at night, that surviving period photographs and tools at the site are consistent with conventional, if labor-intensive, mechanical leverage techniques including tripod hoists, block and tackle systems, and long lever arms rated for exactly this scale of load, and that no antigravity or magnetic levitation mechanism has ever been demonstrated to move stone at all, let alone the specific tonnages Leedskalnin handled, leaving the extraordinary-technique hypothesis with no supporting method where the mundane hypothesis has a complete and physically adequate one.
\end{abstract}

\section{Introduction}

Edward Leedskalnin, a Latvian immigrant, built Coral Castle over nearly three decades, ultimately relocating the entire structure ten miles to its current site in the 1930s, moving multi-ton coral blocks in the process, all reportedly without assistance and without allowing observers to witness his actual working methods, a genuinely unusual and impressive feat of solitary labor that has, since his death, accumulated a substantial body of claims attributing the construction to a discovered antigravity or magnetic levitation principle, sometimes connected to broader claims about how the Egyptian pyramids or other megalithic structures were built.

\section{Why "No One Saw How" Is Not Itself Evidence of an Exotic Method}

Leedskalnin was documented by contemporaries as deliberately private about his working methods, reportedly telling visitors and neighbors that he preferred to work only when unobserved, a personal eccentricity independently attested by multiple period sources and consistent with a solitary immigrant laborer's general reticence rather than requiring an ulterior motive to conceal an exotic technology. The absence of eyewitness testimony to a specific mundane technique is adequately explained by this well-documented preference for privacy, without needing to posit that Leedskalnin's secrecy specifically protected a physics breakthrough rather than simply being, as multiple contemporaries described it, a personal quirk of an unusually solitary man.

\section{What the Site Itself Shows}

Coral Castle preserves period tools and structural evidence, including tripod-like structures, iron implements consistent with wedging and prying techniques, and the coral rock's own natural quarrying properties, coral limestone being relatively soft and workable when freshly cut and hardening with air exposure, that are consistent with established techniques for moving large stone blocks using leverage, rollers, and hoists, methods with a long documented history in vernacular stonemasonry and well within the mechanical advantage achievable by lever arms, fulcrums, and tripod hoists of the kind period photographs of the site show present. Structural engineers and stonemasons who have examined the site, including a 1986 television documentary segment featuring a mason attempting to replicate specific Coral Castle features using period-appropriate hand tools and simple hoisting equipment, have found the required forces and mechanical advantages entirely achievable by a strong, patient individual working alone over the extended multi-decade timescale Leedskalnin actually used, a timescale the antigravity narrative tends to compress or ignore in favor of emphasizing single-block tonnage in isolation from the years available to move each one.

\section{Why the Alternative Hypothesis Has No Demonstrated Method at All}

The antigravity or magnetic levitation narrative attributes Leedskalnin's results to a discovered but never demonstrated or documented physical principle, with no surviving device, no technical description in Leedskalnin's own extensive published pamphlets on magnetism, which discuss his own idiosyncratic and scientifically unconventional theories of magnetic currents but describe no stone-lifting mechanism, and no independent replication by anyone using the claimed principle to move comparable loads in the seven decades since his death. This leaves the exotic hypothesis without a specified mechanism to evaluate at all, while the mundane hypothesis, conventional leverage and hoisting over a multi-decade timeline, is not only specified but independently demonstrated as physically adequate by direct replication attempts of individual site features.

\section{The Entropic Gravity Framing}

Coral Castle's antigravity attribution proposes no specific physical mechanism, resting instead on the general absence of a witnessed construction method, and is included in this series primarily as a case study in how an evidentiary gap, no one saw the technique, gets filled by an appealingly mysterious explanation in the absence of the specific technique's documentation, even when the mundane explanation is independently demonstrated adequate and the exotic explanation is not demonstrated at all.

\section{Objections}

It is sometimes argued that Leedskalnin's own theoretical writings on magnetism, however unconventional, at minimum establish that he believed himself to have discovered genuine and significant principles of magnetic physics, lending some credibility to the idea that his construction technique drew on this belief system. Leedskalnin's surviving magnetism pamphlets are widely regarded by physicists who have reviewed them as containing significant misunderstandings of conventional electromagnetic theory rather than a coherent alternative framework, and more importantly, contain no description, however unconventional, of a stone-lifting application of his ideas, meaning even taking his own theoretical beliefs at face value does not supply the missing mechanism the antigravity construction narrative requires.

\section{Conclusion}

Leedskalnin's solitary, decades-long construction of Coral Castle is fully explained by well-documented personal secrecy combined with conventional leverage and hoisting techniques independently demonstrated adequate to the task, while the antigravity or magnetic levitation alternative offers no demonstrated mechanism, no technical documentation even in Leedskalnin's own extensive personal writings, and no independent replication in the decades since, leaving the mundane explanation as the only one actually supported by evidence.

\end{document}
